\section{总结与下一步工作}
\label{sec:conclusion}

本任务已经形成从数据解析、基础潮流到临界稳定分析的工作流。
NR 具有近解区域的快速收敛能力，FDLF 具有较低的单轮计算成本，但对高 $R/X$ 和强耦合工况更敏感。
接近稳定极限后，算法应当结合最终残差、无功限值、雅可比奇异性和连续潮流分支综合判断。

\subsection{当前结论}

\begin{enumerate}
  \item 常规工况下 NR 与 FDLF 均能可靠求解，前者迭代少，后者单轮成本低；
  \item 全局化方法在本0~300的小规模系统中对于负荷统一增长的病态收敛优化不明显；
  \item 各种参数化的 CPF 各有优劣。与奇异性指标及矩阵条件数可作为最大负荷能力判断的依据。
\end{enumerate}

\subsection{当前不足}

目前 CPF 结果有待进一步整理分析，对于错误分支的产生与特点认识不足。
规模分析也仍基于当前 Python 与矩阵实现，尚未区分算法复杂度。

\subsection{下一步工作}

后续工作建议如下：
\begin{enumerate}
  \item 统一不同求解器的负荷参数化、PV--PQ 转换与残差口径；
  \item 增加雅可比最小奇异值、条件数和临界特征向量监测；
  \item 引入稀疏矩阵存储与稀疏 LU/QR 分解；
  \item 在故障、非均匀负荷增长等场景下复核 CPF 参数化选择，并区分纯参数化与后备接管的贡献；
  \item 保存配置、代码版本、原始数据和图像来源，提高结果可复现性。
\end{enumerate}


\subsection{探索性内容}

按照《高等电力网络分析 (张伯明，严正著)》中的证明
及《PQ分解法解潮流方程收敛性的影响因素分析_王若涵》中关于迭代顺序的研究进行讨论验证。